
Determining the electronic structures is probably the most important rationale behind the application of the electronic structure methods (DFT and GW) in materials modelling. Fundamentally, it enables the characterisations of materials into metals, semiconductors and insulators as their basic electronic properties, which dictate their practical applications and many other associated physical and chemical properties. For examples, by projecting the electronic eigenfunctions to atomic orbital basis functions, one can obtain further fundamental insights into the origins of electronic states near the Fermi level from the perspective of orbital hybridisations,\cite{dronskowski1993crystal} which also play important roles in dictating the materials' stabilities. For instance, in lead halide perovskites, the distortion of atomic positions along particular phonon modes can lead to redistribution of the electronic charge densities to increase its structural stability,\cite{yang2020composition} which can be viewed as a manifestation of the pseudo-Jahn-Teller effect.\cite{bersuker2013pseudo} The topological information of the electronic band structure can be utilised directly to compute the electronic conductivity tensors for a material subjected to external fields (electric, magnetic, thermal gradients, and so on). The Kohn-Sham eigenenergies and eigenfunctions serve as the critical ingredient for computing the frequency-dependent dielectric function  [\cref{eq:adler-wiser}], which is not only an important step in GW calculation, but also useful for computing the optical spectra of semiconductors.\cite{ramasubramaniam2012large} Finally, the band structure is also an important component in facilitating the understanding of more complex many-body phenomena, such as band-mediated magnetisms,\cite{walsh2008theoretical} in solid state materials.

Formally, the electronic band gap energy $E_g$ of a system should be taken as the energy difference between the LUMO and HOMO states, i.e. $E_g=\varepsilon_{\mbox{\scriptsize{LUMO}}}-\varepsilon_{\mbox{\scriptsize{HOMO}}}$, which, as can be seen in \cref{eq:ionisation_potential} and \cref{eq:electron_affinity}, that, fundamentally, $E_g$ is determined not only by the ground, but also the excited state energies of a many-electron system. However, DFT calculation only gives the ground-state energy of an $N$ electron system with fixed orbital occupation under the mean-field approximation. What enables us to overcome this limitation, and extract the band gap energies from a ground-state DFT calculation is the Janak theorem,\cite{janak1978proof} which proved that $\partial E/\partial c_i=\varepsilon_i$, $i.e.$ the derivative of the Kohn-Sham total energy with respect to the occupation of the $i$-th eigenorbital is the corresponding orbital eigenenergy. The theory was established based on a generalisation of the Kohn-Sham theory, whereby in the electron density $\rho(\bm{r})=\sum_i c_i|\varphi_i(\bm{r})|^2$ [\cref{eq:electron_density_KS_orbital}], one considered the case where the orbital occupation $c_i$ was allowed to be a continuous (fractional) number between 0 and 2 (constraint of the Pauli exclusion principle). With this, the orbital eigenenergy can be written as a function of the orbital occupancy, or $\varepsilon(n_i)$, and $E(N+1)-E(N)=\int_0^1 \varepsilon(n_i)dn_i$. This serves as the fundamental basis to compute the electronic band structures of materials with DFT. Nevertheless, one should note that the problem of mean-field DFT in underestimating $E_{g}$ (\cref{sect:DFA_challenge}) does not invalidate the Janak's theorem, but a problem of approximating the many-body exchange-correlation potential with an one-electron density matrix. The application of post-DFT method, such as the GW approximation, in which the quasiparticle energy now follows the definition of  \cref{eq:linearised_QP_eqn}, should (partially) overcome the limitations in mean-field DFT, leading to more accurate estimation of the electronic energy levels.

Practically, when using the planewave basis functions [\cref{eq:bloch_wave}] to solve the Kohn-Sham equation for solids, one first seeks to construct the basis function on a series of discrete $k$-points in the reciprocal space, either along high-symmetry $k$-point paths for calculating the band structures or evenly spaced within the first Brillouin zone (BZ). At each selected $k$-point, solving the Kohn-Sham equation will give rise to a spectrum of energy eigenvalues $\{\varepsilon_{i\bm{k}}\}$, in which $i$ is often referred to as the band index. The electronic band structure can thus be constructed by connecting $\varepsilon_{i\bm{k}}$ that belong to the same band index with different $k$ by following the  $k$-point paths that were pre-defined based on the crystal symmetry, whereas the electronic density-states (DOS) as a function of electronic energy levels are computed by integrating across all $k$-points in the first BZ as:
\begin{equation}
    \label{eq:electronic_DOS}
    \rho(\varepsilon)=\sum_i\int_{\mbox{\scriptsize{BZ}}} \frac{\mbox{d}\bm{k}}{\Omega_{\mbox{\scriptsize{BZ}}}}\delta(\varepsilon-\varepsilon_{i\bm{k}}).
\end{equation}
